% Options for packages loaded elsewhere
\PassOptionsToPackage{unicode}{hyperref}
\PassOptionsToPackage{hyphens}{url}
\documentclass[
]{article}
\usepackage{xcolor}
\usepackage{amsmath,amssymb}
\setcounter{secnumdepth}{-\maxdimen} % remove section numbering
\usepackage{iftex}
\ifPDFTeX
  \usepackage[T1]{fontenc}
  \usepackage[utf8]{inputenc}
  \usepackage{textcomp} % provide euro and other symbols
\else % if luatex or xetex
  \usepackage{unicode-math} % this also loads fontspec
  \defaultfontfeatures{Scale=MatchLowercase}
  \defaultfontfeatures[\rmfamily]{Ligatures=TeX,Scale=1}
\fi
\usepackage{lmodern}
\ifPDFTeX\else
  % xetex/luatex font selection
\fi
% Use upquote if available, for straight quotes in verbatim environments
\IfFileExists{upquote.sty}{\usepackage{upquote}}{}
\IfFileExists{microtype.sty}{% use microtype if available
  \usepackage[]{microtype}
  \UseMicrotypeSet[protrusion]{basicmath} % disable protrusion for tt fonts
}{}
\makeatletter
\@ifundefined{KOMAClassName}{% if non-KOMA class
  \IfFileExists{parskip.sty}{%
    \usepackage{parskip}
  }{% else
    \setlength{\parindent}{0pt}
    \setlength{\parskip}{6pt plus 2pt minus 1pt}}
}{% if KOMA class
  \KOMAoptions{parskip=half}}
\makeatother
\usepackage{color}
\usepackage{fancyvrb}
\newcommand{\VerbBar}{|}
\newcommand{\VERB}{\Verb[commandchars=\\\{\}]}
\DefineVerbatimEnvironment{Highlighting}{Verbatim}{commandchars=\\\{\}}
% Add ',fontsize=\small' for more characters per line
\newenvironment{Shaded}{}{}
\newcommand{\AlertTok}[1]{\textcolor[rgb]{1.00,0.00,0.00}{\textbf{#1}}}
\newcommand{\AnnotationTok}[1]{\textcolor[rgb]{0.38,0.63,0.69}{\textbf{\textit{#1}}}}
\newcommand{\AttributeTok}[1]{\textcolor[rgb]{0.49,0.56,0.16}{#1}}
\newcommand{\BaseNTok}[1]{\textcolor[rgb]{0.25,0.63,0.44}{#1}}
\newcommand{\BuiltInTok}[1]{\textcolor[rgb]{0.00,0.50,0.00}{#1}}
\newcommand{\CharTok}[1]{\textcolor[rgb]{0.25,0.44,0.63}{#1}}
\newcommand{\CommentTok}[1]{\textcolor[rgb]{0.38,0.63,0.69}{\textit{#1}}}
\newcommand{\CommentVarTok}[1]{\textcolor[rgb]{0.38,0.63,0.69}{\textbf{\textit{#1}}}}
\newcommand{\ConstantTok}[1]{\textcolor[rgb]{0.53,0.00,0.00}{#1}}
\newcommand{\ControlFlowTok}[1]{\textcolor[rgb]{0.00,0.44,0.13}{\textbf{#1}}}
\newcommand{\DataTypeTok}[1]{\textcolor[rgb]{0.56,0.13,0.00}{#1}}
\newcommand{\DecValTok}[1]{\textcolor[rgb]{0.25,0.63,0.44}{#1}}
\newcommand{\DocumentationTok}[1]{\textcolor[rgb]{0.73,0.13,0.13}{\textit{#1}}}
\newcommand{\ErrorTok}[1]{\textcolor[rgb]{1.00,0.00,0.00}{\textbf{#1}}}
\newcommand{\ExtensionTok}[1]{#1}
\newcommand{\FloatTok}[1]{\textcolor[rgb]{0.25,0.63,0.44}{#1}}
\newcommand{\FunctionTok}[1]{\textcolor[rgb]{0.02,0.16,0.49}{#1}}
\newcommand{\ImportTok}[1]{\textcolor[rgb]{0.00,0.50,0.00}{\textbf{#1}}}
\newcommand{\InformationTok}[1]{\textcolor[rgb]{0.38,0.63,0.69}{\textbf{\textit{#1}}}}
\newcommand{\KeywordTok}[1]{\textcolor[rgb]{0.00,0.44,0.13}{\textbf{#1}}}
\newcommand{\NormalTok}[1]{#1}
\newcommand{\OperatorTok}[1]{\textcolor[rgb]{0.40,0.40,0.40}{#1}}
\newcommand{\OtherTok}[1]{\textcolor[rgb]{0.00,0.44,0.13}{#1}}
\newcommand{\PreprocessorTok}[1]{\textcolor[rgb]{0.74,0.48,0.00}{#1}}
\newcommand{\RegionMarkerTok}[1]{#1}
\newcommand{\SpecialCharTok}[1]{\textcolor[rgb]{0.25,0.44,0.63}{#1}}
\newcommand{\SpecialStringTok}[1]{\textcolor[rgb]{0.73,0.40,0.53}{#1}}
\newcommand{\StringTok}[1]{\textcolor[rgb]{0.25,0.44,0.63}{#1}}
\newcommand{\VariableTok}[1]{\textcolor[rgb]{0.10,0.09,0.49}{#1}}
\newcommand{\VerbatimStringTok}[1]{\textcolor[rgb]{0.25,0.44,0.63}{#1}}
\newcommand{\WarningTok}[1]{\textcolor[rgb]{0.38,0.63,0.69}{\textbf{\textit{#1}}}}
\usepackage{longtable,booktabs,array}
\usepackage{caption}
\captionsetup[table]{skip=6pt}
\newcounter{none} % for unnumbered tables
\usepackage{calc} % for calculating minipage widths
% Correct order of tables after \paragraph or \subparagraph
\usepackage{etoolbox}
\makeatletter
\patchcmd\longtable{\par}{\if@noskipsec\mbox{}\fi\par}{}{}
\makeatother
\setlength{\emergencystretch}{3em} % prevent overfull lines
\providecommand{\tightlist}{%
  \setlength{\itemsep}{0pt}\setlength{\parskip}{0pt}}
\usepackage{bookmark}
\IfFileExists{xurl.sty}{\usepackage{xurl}}{} % add URL line breaks if available
\urlstyle{same}
% fallback for those not using the hyperref driver hyperxmp:
\makeatletter
\@ifundefined{xmpquote}{\newcommand{\xmpquote}[1]{#1}}{}
\makeatother
\hypersetup{
  pdftitle={Retaining a metric basis under one sensor erasure},
  pdfauthor={timofy (project coordinator)},
  hidelinks,
  pdfcreator={LaTeX via pandoc}}

\title{Retaining a metric basis under one sensor erasure}
\author{timofy (project coordinator)}
\date{}

\begin{document}
\maketitle

\section{Research note}\label{research-note}

AI-assisted research draft. Methodology, contributions and review status
are described in Section 11 and the companion methodology note.

\section{Abstract}\label{abstract}

For a connected graph G and a minimum resolving set B, we consider the
minimum cardinality rho(G,B) of a fault-tolerant resolving set
containing B, and compare it with the unconstrained fault-tolerant
metric dimension beta2(G). The difference p(G,B)=rho(G,B)-beta2(G)
measures a cardinality penalty for retaining the original base. Using
established ordinary-resolution characterizations for unicyclic graphs,
we prove that p(G,B)\textless=1 for every metric basis of every
connected simple unweighted subcubic unicyclic graph, and give a
seven-vertex example attaining equality. An accompanying exact
enumeration through ten vertices checks 331 graph classes and all 3,900
bases: 3,761 bases have penalty zero and 139 have penalty one. A
separately authored computational audit verifies coverage, all optima
and 237,928 subset certificate entries. This is a narrow structural
result in pure graph theory; it does not establish an efficient
optimization algorithm or a physical sensor-deployment improvement.

\section{1. Problem and literature
context}\label{problem-and-literature-context}

Ordinary metric dimension minimizes the number of vertices needed to
distinguish all vertices by their distance vectors. Fault-tolerant
metric dimension requires resilience to erasure of a sensor coordinate.
We study the additional cost of retaining an arbitrary prescribed metric
basis when upgrading a unicyclic graph to one-fault-tolerant resolution.
For connected subcubic unicyclic graphs, this cost relative to an
unconstrained optimal fault-tolerant resolving set is at most one, and
the bound is sharp.

Preserving a given resolving set during augmentation is established
prior art. Hernando, Mora, Slater and Wood {[}1 Lemma 3.1{]} start with
a given resolving set and construct a fault-tolerant superset using
closed neighborhoods and additional neighborhood-containment witnesses.
Their Theorem 3.3 gives a bound in terms of ordinary metric dimension.
Their Remark 2.4 allows arbitrary choices of two vertices in the two
terminal paths of a tree stem with exactly two terminal paths, preceding
the local freedom used in our terminal-fork exchange. Our comparison
concerns the \emph{optimal} superset of a specified basis relative to
the unconstrained fault-tolerant optimum, rather than claiming
augmentation itself as new.

Sedlar and Skrekovski {[}2{]} provide the ordinary unicyclic
resolving-set characterization essential to the proof. Their Lemmas 1,
3, 6 and 7 supply branch coverage, biactivity, geodesic triples and the
A/B/C obstruction characterization; Theorem 11 gives ordinary dimension.
Our exchange lemma is an application/corollary of that characterization,
with the tree predecessor just noted. It is not a new characterization
or a separate major novelty claim. The one-branch-root repair and
antipodal completion supply additional arguments for the
prescribed-basis comparison.

Knor, Sedlar and Skrekovski {[}3{]} study how expensive fault tolerance
can be relative to ordinary metric dimension, including exponential
growth and higher multiplicities. Although the language includes
upgrading a metric basis, the optimized invariant is unconstrained
fault-tolerant dimension. Geneson and Tsai {[}4 Section 2 and the proof
of Theorem 2.7{]} explicitly extend arbitrary resolving sets using their
closed neighborhoods and further witnesses, another predecessor for
basis-preserving augmentation. Their stated bounds and extremal
questions concern ordinary dimension. The inspected arXiv version has
differing constants in its abstract and Theorem 2.7; no exact numerical
bound from it is used here.

Related terminology has different quantifiers. Hakanen et al.~{[}5{]}
study vertices contained in every or in no minimum ordinary metric
basis, rather than vertices imposed as part of an input. Their Section
3.1.1 also records the known dimension-two result for odd type-1
unicyclic graphs and dimension-two-or-three result for even ones,
attributing these to Poisson and Zhang (2002). Our branch-free dimension
facts are therefore acknowledged as known; the explicit odd-cycle
calculation below is included for completeness. Bergougnoux, Defrain and
Mc Inerney {[}6 Sections 6-7{]} use extension problems with required and
forbidden vertices in enumeration of inclusion-minimal solutions.
Ordinary inclusion-minimal extension differs from minimum-cardinality
fault-tolerant augmentation of an already minimum basis.

Estrada-Moreno {[}7{]} studies k-metric dimension; k=2 is the
unconstrained fault-tolerant invariant here. Theorems 4-5 concern type-1
graphs with exactly one degree-three cycle vertex; Theorem 6 concerns at
least three special branching vertices under further hypotheses. These
do not state an extension guarantee for every prescribed minimum basis.
Our seven-vertex witness has two degree-three cycle vertices, so it is
not the single-pendant-path case. Prabhu et al.~{[}8 version 2{]} use
attaching fault-tolerant sets chosen outside an attachment set: their
additional selected vertices are subject to deletion while the
designated attachments remain available. Here every sensor, including
every member of B, is subject to erasure. Their point-attachment
formulas also have structural conditions.

We did not find this prescribed-basis comparison in the sources checked.
The result-specific public-source search, completed on 18 October 2026,
judged the sharp theorem apparently new with moderate confidence, not an
originality guarantee. The approximately ten-minute search covered
English-language indexed literature and adjacent references, including
fixed/prescribed/forced landmarks, extensions, supersets, exchange and
disjoint resolving sets. A second terminology pass was by the same
searcher, not an independent second literature review. It was not an
exhaustive MathSciNet/zbMATH search and incompletely covered theses,
non-English and unindexed work or unavailable texts. A potentially
relevant pseudotree survey was identified but not retrieved and is not
used to certify absence of prior art. No authors were contacted. No
separate priority claim is made for the elementary union bound, the
standard graph family or the technical exchange lemma.

\section{2. Definitions and
conventions}\label{definitions-and-conventions}

Graphs in the original pilot are finite, connected, simple, undirected
and unweighted. Let d\_G(u,v) be the number of edges in a shortest u-v
path. For distinct u,v define

\[D_G(u,v)=\{s\in V(G):d_G(u,s)\ne d_G(v,s)\}.\]

A set S resolves G when S intersects every D\_G(u,v). Its minimum
possible size is beta(G). We call a resolving set B a metric basis only
when \textbar B\textbar=beta(G); mere inclusion-minimality is
insufficient.

A set S is fault-tolerant when \textbar S intersect
D\_G(u,v)\textbar\textgreater=2 for every distinct pair u,v.
Equivalently, S without s resolves G for each s in S, with the empty set
excluded on the graphs considered here. To see equivalence, two
distinguishers leave at least one after any erasure. If a pair has at
most one distinguisher, either it is already unresolved or erasure of
that sole distinguisher makes it unresolved. Distances are always taken
in the unchanged graph G. Erasing a sensor does not delete its vertex,
incident edges or paths through it. The failed sensor vertex remains in
the target population and distances are not recomputed.

Write FT(G) for the family of fault-tolerant sets and define

\[\beta_2(G)=\min_{S\in\operatorname{FT}(G)}|S|,\qquad
\rho(G,B)=\min_{S\in\operatorname{FT}(G),\,B\subseteq S}|S|,
\qquad p(G,B)=\rho(G,B)-\beta_2(G).\]

The subscript in beta2 denotes the two-distinguishers condition; it is
not an assertion that every cited source uses identical notation. All
optima exist: for every pair u,v, both u and v distinguish that pair, so
V(G) is fault-tolerant. Containment restricts feasible sets, hence
p(G,B)\textgreater=0. The quantity rho counts all sensors in the
superset; rho-\textbar B\textbar{} counts additions. These are different
statistics.

A connected graph is unicyclic when it has exactly one cycle; for a
simple connected graph this is equivalent to
\textbar E\textbar=\textbar V\textbar. Subcubic means maximum degree at
most three. The finite domain is
3\textless=\textbar V\textbar\textless=10; no simple unicyclic graph has
fewer than three vertices. Pure cycles are included.

\section{3. Audited finite result}\label{audited-finite-result}

\textbf{Computational Result 1 (finite domain).} For every connected
simple unweighted subcubic unicyclic graph G with
3\textless=\textbar V(G)\textbar\textless=10, and for every metric basis
B of G, p(G,B)\textless=1. In this domain both values zero and one
occur.

This statement is supported by exhaustive integer computation, the
coverage arguments in Section 4 and an independent executable audit. It
is not a proof-assistant-checked theorem and is not asserted for orders
above ten.

{\def\LTcaptype{none} % do not increment counter
\begin{longtable}[]{@{}
  >{\raggedleft\arraybackslash}p{(\linewidth - 10\tabcolsep) * \real{0.1667}}
  >{\raggedleft\arraybackslash}p{(\linewidth - 10\tabcolsep) * \real{0.1667}}
  >{\raggedleft\arraybackslash}p{(\linewidth - 10\tabcolsep) * \real{0.1667}}
  >{\raggedleft\arraybackslash}p{(\linewidth - 10\tabcolsep) * \real{0.1667}}
  >{\raggedleft\arraybackslash}p{(\linewidth - 10\tabcolsep) * \real{0.1667}}
  >{\raggedleft\arraybackslash}p{(\linewidth - 10\tabcolsep) * \real{0.1667}}@{}}
\toprule\noalign{}
\begin{minipage}[b]{\linewidth}\raggedleft
Order n
\end{minipage} & \begin{minipage}[b]{\linewidth}\raggedleft
Graph classes
\end{minipage} & \begin{minipage}[b]{\linewidth}\raggedleft
Metric bases
\end{minipage} & \begin{minipage}[b]{\linewidth}\raggedleft
Bases with p=0
\end{minipage} & \begin{minipage}[b]{\linewidth}\raggedleft
Bases with p=1
\end{minipage} & \begin{minipage}[b]{\linewidth}\raggedleft
Graphs admitting p=1
\end{minipage} \\
\midrule\noalign{}
\endhead
\bottomrule\noalign{}
\endlastfoot
3 & 1 & 3 & 3 & 0 & 0 \\
4 & 2 & 9 & 9 & 0 & 0 \\
5 & 4 & 29 & 29 & 0 & 0 \\
6 & 9 & 74 & 74 & 0 & 0 \\
7 & 17 & 175 & 170 & 5 & 2 \\
8 & 38 & 407 & 394 & 13 & 5 \\
9 & 80 & 971 & 943 & 28 & 11 \\
10 & 180 & 2,232 & 2,139 & 93 & 27 \\
Total & 331 & 3,900 & 3,761 & 139 & 45 \\
\end{longtable}
}

These are unweighted counts of bases on one labeled representative per
isomorphism class, not counts over all n! relabelings and not counts
modulo automorphisms of a representative. Every tied minimum base is
retained. The displayed relative frequencies therefore describe this
counting convention rather than a random-graph probability model.

\section{4. Complete generation and equivalence
reduction}\label{complete-generation-and-equivalence-reduction}

\subsection{4.1. Spanning-tree
construction}\label{spanning-tree-construction}

\textbf{Proposition 2 (coverage).} Every subcubic unicyclic graph on n
vertices is produced by taking a subcubic tree on n vertices and adding
a non-edge whose endpoints have tree degree at most two.

\emph{Proof.} Delete an edge of the graph's unique cycle. The resulting
graph is a spanning tree and remains subcubic. The endpoints of the
deleted edge each lose one incident edge and therefore have degree at
most two. The original graph is recovered by adding that non-edge.
Conversely, adding a non-edge to a tree creates one cycle and preserves
connectivity; the endpoint restriction preserves maximum degree three.
The graph is simple because an existing edge is never added. QED.

All subcubic trees are generated inductively by adding a new leaf at
each vertex of degree at most two, starting with one vertex. Deleting
any leaf of any nontrivial subcubic tree yields a smaller subcubic tree
and makes the original attachment vertex eligible for reattachment.
Induction therefore establishes coverage. Deduplication is applied at
every order, but no structural subclass is sampled or omitted.

\subsection{4.2. Exact isomorphism keys}\label{exact-isomorphism-keys}

A rooted tree is encoded recursively by a matched parenthesis pair
enclosing its sorted child encodings. Induction on size shows that two
encodings agree exactly when the rooted trees are isomorphic: the sorted
child lists agree precisely as multisets of rooted isomorphism classes.
For an unrooted tree the primary implementation takes the minimum such
encoding over all choices of root. Equality gives an isomorphism for a
pair of roots, while an unrooted isomorphism induces a bijection of the
root choices.

For a unicyclic graph, repeatedly peeling leaves identifies its unique
cycle. At each cycle vertex, encode the rooted off-cycle attachment,
excluding the cycle edges. Form the sequence of these rooted codes in
cycle order. The canonical key is the least sequence over all cyclic
rotations and both orientations, with explicit separators and cycle
length. An isomorphism must map the unique cycle to the unique cycle and
acts on it by a dihedral symmetry; it also induces rooted isomorphisms
of the attachments. Conversely, matching sequences and attachment
isomorphisms assemble to a graph isomorphism. Thus these full strings,
rather than hash collisions or numerical invariants, decide equality of
primary graph classes.

\subsection{4.3. Separate coverage
construction}\label{separate-coverage-construction}

A second construction chooses cycle length k, distributes n-k vertices
among its k positions, and attaches either nothing or one rooted tree at
each position. Every off-cycle vertex has at most two children; each
cycle vertex has at most one off-cycle child. Recursive generation of
unordered zero-, one- and two-child rooted trees covers all such
attachments. Cutting the unique cycle edges of any target graph shows
that every target has this form.

The same-task verifier deduplicated and matched these constructions
through exact degree-preserving adjacency backtracking, independently of
the primary canonical equality routine. The separate external auditor
independently implemented the cycle/attachment construction and
canonical encoding and obtained the same class set at every order. The
external construction is independently authored; it is not claimed to be
a third distinct structural characterization beyond the tree-addition
and cycle-attachment routes.

\subsection{4.4. Why representatives suffice for all labels and
bases}\label{why-representatives-suffice-for-all-labels-and-bases}

An isomorphism f preserves distances. It therefore carries resolving
sets and fault-tolerant sets bijectively to their counterparts,
preserves cardinalities and preserves containment B subseteq S.
Consequently beta, beta2, rho and p are invariant under simultaneous
transport of G and B. Since every minimum base of each chosen
representative is explicitly checked, every base on every relabeling is
covered. No reduction to a single base or a single favorable tie is
used.

\section{5. Exact optimization and
certificates}\label{exact-optimization-and-certificates}

The primary implementation computes shortest-path distances by
breadth-first search. Vertices are labeled 0,\ldots,n-1; bit i of a mask
denotes vertex i. For each unordered pair it stores the integer mask of
D\_G(u,v). It visits all 2\^{}n subsets and tests whether every pair has
at least one or at least two selected distinguishers. Minimum
cardinalities give beta and beta2, and all resolving masks at
cardinality beta are retained. For each such B it scans the
fault-tolerant masks containing B and minimizes their total cardinality.
Tied optimum witnesses are chosen by increasing integer mask; choosing
one witness does not discard any minimum base.

Each graph record contains its edge list, distance matrix, beta, beta2
and an optimum fault-tolerant witness. Each base record contains B, rho,
p and an optimum superset. The complete subset certificate records its
cardinality and, when a condition fails, a concrete vertex pair
witnessing failure. These records certify lower bounds through
exhaustive failed smaller subsets as well as upper bounds through
feasible sets. At n\textless=10 there are at most 1,024 subsets and 45
pairs per graph. The direct base/superset comparison is bounded by
4\^{}n comparisons per graph. All arithmetic relevant to the mathematics
is integer arithmetic.

The same-task verifier instead uses Floyd-Warshall distances and tuple
combinations. Resolving is tested by uniqueness of vertex distance
signatures; fault tolerance is tested by deletion of each sensor
coordinate. A subsequent serialization audit parses the saved outputs
and repeats verification. It also checks 1,182 deterministic relabelings
for canonical-key invariance and transport of witness
containment/cardinality. These relabeling tests supplement, rather than
replace, the exact isomorphism argument.

The external auditor wrote its core before examining the original
implementation or saved result contents. It imports no pilot module. It
computes distances by BFS, constructs all subset signature tables and
tests one-coordinate deletion. It checks every recorded optimum, the
complete base families, all 237,928 subset entries, distance matrices
and all supplied obstruction pairs. Its audit does not depend on the
convention that an optimum witness or failing pair is lexicographically
first. The external verdict is PASS for the finite domain. The source
hashes remained unchanged during its successful run.

\section{6. A seven-vertex example with a direct
proof}\label{a-seven-vertex-example-with-a-direct-proof}

Let V=\{0,1,2,3,4,5,6\} and

\[E=\{01,02,03,14,25,45,56\}.\]

Here ij denotes the undirected edge \{i,j\}. The cycle is 0-1-4-5-2-0;
leaves 3 and 6 attach to 0 and 5. Let B=\{0,5\}.

\textbf{Proposition 3.} This graph and base satisfy beta(G)=2,
beta2(G)=3 and rho(G,B)=4, hence p(G,B)=1.

\emph{Proof.} The following table gives signatures in the indicated
sensor order.

{\def\LTcaptype{none} % do not increment counter
\begin{longtable}[]{@{}rlll@{}}
\toprule\noalign{}
Vertex & B=(0,5) & C=(1,2) & T=(1,2,3) \\
\midrule\noalign{}
\endhead
\bottomrule\noalign{}
\endlastfoot
0 & (0,2) & (1,1) & (1,1,1) \\
1 & (1,2) & (0,2) & (0,2,2) \\
2 & (1,1) & (2,0) & (2,0,2) \\
3 & (1,3) & (2,2) & (2,2,0) \\
4 & (2,1) & (1,2) & (1,2,3) \\
5 & (2,0) & (2,1) & (2,1,3) \\
6 & (3,1) & (3,2) & (3,2,4) \\
\end{longtable}
}

Every B-signature is distinct, so B resolves. No singleton resolves: for
singleton sensors 0,1,2,3,4,5,6 respectively, the equal-distance pairs
\{1,2\}, \{0,4\}, \{0,5\}, \{1,2\}, \{0,2\}, \{0,1\}, \{0,1\} are
obstructions. Hence beta(G)=2 and B is a metric basis.

Every two distinct T-signatures differ in at least two coordinates, as
direct inspection of the seven displayed triples verifies. Thus T is
fault-tolerant and beta2\textless=3. If a set of cardinality at most two
were fault-tolerant, erasing a sensor would leave at most a singleton
resolving set, impossible here. Hence beta2=3.

There are five possible size-three supersets of B. The following table
gives a pair with only one distinguishing sensor in each, together with
that sensor. Its erasure causes the displayed collision.

{\def\LTcaptype{none} % do not increment counter
\begin{longtable}[]{@{}rlr@{}}
\toprule\noalign{}
Added vertex a & Pair & Sole distinguishing sensor in B union \{a\} \\
\midrule\noalign{}
\endhead
\bottomrule\noalign{}
\endlastfoot
1 & \{2,3\} & 5 \\
2 & \{1,3\} & 5 \\
3 & \{1,2\} & 5 \\
4 & \{2,6\} & 0 \\
6 & \{0,1\} & 0 \\
\end{longtable}
}

None is fault-tolerant, so rho(G,B)\textgreater=4. The C-signatures are
also distinct, and B and C are disjoint. Their union is fault-tolerant:
after erasure of a vertex from one of the two sets, the other resolving
set remains intact. Thus B union C=\{0,1,2,5\} is a feasible size-four
superset and rho(G,B)=4. QED.

For completeness, the saved exhaustive certificate for this graph
contains all 128 subsets and all 15 metric bases. It is additional
evidence, not needed for the displayed direct argument. Order seven is
the first order admitting positive penalty in the completed finite
enumeration; this does not assert uniqueness of the example.

\section{7. General bound}\label{general-bound}

\textbf{Proof dependencies and review.} The complete argument uses the
cited ordinary-resolution characterization. Separate AI sessions
logically reviewed the proof, including a referee verdict of ``proved as
written'' for the recorded proof revision. This is not human or journal
peer review. The integrated elementary size-bound version and the
reviewed formula-based revision are distinguished in the accompanying
review report. The finite result in Section 3 stands independently and
is not used to infer the arbitrary-order bound.

\textbf{Theorem 4.} Let G be a connected simple unweighted subcubic
unicyclic graph, and let B be any metric basis of G. Then

\[0\le \rho(G,B)-\beta_2(G)\le 1.\]

The upper bound is sharp, as shown by Proposition 3. The proof follows
from the lemmas and cases below, using the explicitly identified
ordinary-resolution results of Sedlar and Skrekovski.

For the propositions, extend rho(G,B) and p(G,B) by the same formulas to
any nonempty resolving set B, whether or not it is minimum. When
discussing metric dimension we explicitly return to a metric basis.

\subsection{7.1. A general union bound}\label{a-general-union-bound}

\textbf{Lemma 5 (union bound).} If B is a nonempty resolving set of a
connected graph with at least two vertices, then
p(G,B)\textless=\textbar B\textbar-1. In particular, every metric basis
of size two satisfies p\textless=1.

\emph{Proof.} Let F be a minimum fault-tolerant set. If F intersects B,
the fault-tolerant superset F union B has cardinality at most
\textbar F\textbar+\textbar B\textbar-1. If the sets are disjoint,
choose any f in F. Both F without f and B resolve, and are disjoint.
Their union is fault-tolerant: erasing any one sensor leaves at least
one of the two resolving sets intact. It contains B and has
\textbar F\textbar+\textbar B\textbar-1 elements. Subtracting
beta2=\textbar F\textbar{} gives the bound. QED.

\subsection{7.2. Branch definitions and the imported
characterization}\label{branch-definitions-and-the-imported-characterization}

Let C be the unique cycle. Removing the cycle edges partitions G into
rooted trees T\_v, one at each v in C. In a subcubic graph a
\emph{terminal fork} w is an off-cycle degree-three vertex with two
child paths ending at leaves and with only degree-two internal vertices.
These two paths are its terminal leafward threads. Let U\_w be the union
of their vertices, excluding w. Distinct U\_w are disjoint. Let L be the
number of terminal forks. A \emph{branch-active root} is a cycle root v
whose T\_v contains an off-cycle degree-three vertex; let b count these
roots. This graph-dependent notion is distinct from a root being
\emph{active for S}, meaning that S intersects T\_v. A \emph{direct
cycle thread} is a leafward path attached at a cycle root with only
degree-two internal vertices; its occupancy records whether S meets the
path away from the root.

The dictionary to the cited source is especially simple here. A cycle
vertex has at most one pendant thread. An off-cycle degree-three vertex
has a direction leading toward the cycle and contributes a nonvacuous
thread-coverage requirement exactly when its other two directions are
terminal threads. Each such fork contributes one to the source parameter
L(G), so L(G)=L. Being branch-resolving is therefore precisely the
condition S intersect U\_w is nonempty for every terminal fork w. This
identifies both the source's branch parameter and its branch-resolving
hypothesis before their use below.

We use Sedlar and Skrekovski {[}2{]}, with numbering fixed to
arXiv:2104.00577v1. Lemma 1 states that resolving sets are biactive and
branch-resolving. Definition 5 and Lemmas 6-7 characterize resolution
among such sets by excluding configurations A, B and C. These
configurations depend on the active cycle roots, unoccupied pendant
threads attached directly to the cycle, and the fixed thread lengths of
the graph. The characterization applies to arbitrary biactive
branch-resolving sets, not only minimum ones.

Consequently, if R is resolving and R' is branch-resolving with the same
active cycle roots and the same occupancy of direct cycle threads as R,
then R' resolves. The underlying graph and its thread lengths are
unchanged in this use. In the subcubic setting, branch coverage requires
precisely that every U\_w be met: off-cycle vertices have at most two
pendant threads, and cycle vertices at most one. We also use Lemma 3 of
the same source: a branch-resolving set resolves if its active roots
contain a geodesic triple. For cycle length g, three distinct roots
u,v,w form such a triple when d\_C(u,v)+d\_C(v,w)+d\_C(w,u)=g.

These are the source inputs identified and checked against the full
preprint in the coordinating session-supplied proof review. The
following proof uses elementary basis-size bounds rather than the
stronger metric-dimension formula in Theorem 11. It therefore does not
require that formula as an additional external premise.

\subsection{7.3. Terminal-thread
exchange}\label{terminal-thread-exchange}

\textbf{Lemma 6 (terminal-thread exchange).} For a nonempty resolving
set B in a subcubic unicyclic graph,

\[p(G,B)\le |B|-L.\]

\emph{Proof.} At a terminal fork w, let u\_1,u\_2 be the first vertices
on its two leafward threads. A vertex outside U\_w has equal distance to
u\_1 and u\_2, whereas every vertex in U\_w has different distances to
them. Therefore D\_G(u\_1,u\_2)=U\_w. Every resolving set meets U\_w,
and every fault-tolerant set has at least two vertices in U\_w.

Choose B\_0 by selecting one vertex of B in each U\_w, so
\textbar B\_0\textbar=L. Starting with a minimum fault-tolerant F, for
each selected B\_0 vertex not already present replace one F vertex in
the same U\_w by that vertex. Disjointness of the U\_w ensures these
replacements have no interference or duplication. Write F' for the
resulting set and let phi:F-\textgreater F' be the bijection matching
each replaced vertex to its replacement and fixing the others. Then
\textbar F'\textbar=\textbar F\textbar{} and B\_0 is contained in F'.

Fix z in F', and put x=phi\^{}\{-1\}(z). The set F without x resolves.
The set F' without z still meets every U\_w, since F had at least two
vertices in each such set and replacements preserve the count in each
U\_w, including after matched deletion. Replacements occur within the
same attached rooted tree, so active cycle roots are unchanged relative
to F without x. They occur on threads of off-cycle forks, not on direct
cycle threads, so direct cycle thread occupancy is unchanged too. The
imported characterization implies that F' without z resolves. This holds
for every z, hence F' is fault-tolerant.

Adding B without B\_0 to F' gives a fault-tolerant superset of B of
cardinality at most \textbar F\textbar+\textbar B\textbar-L. Subtract
beta2=\textbar F\textbar. QED.

\textbf{Consequence for metric bases.} If beta=L, then p=0 by this bound
and nonnegativity. The elementary size bounds below will give the
remaining case division.

\subsubsection{Elementary basis-size
bounds}\label{elementary-basis-size-bounds}

\textbf{Lemma 7.} If b\textgreater=2, then beta\textless=L+1; if b=1,
then beta\textless=L+2; and if b=0, then beta\textless=3.

\emph{Proof.} First, any two distinct vertices a,b on a cycle of length
g\textgreater=3 can be completed to a geodesic triple. Choose a shortest
a-b arc of length d\textless=floor(g/2); the complementary arc has
length ell=g-d, with 2\textless=ell\textless=g-1. Choose an interior
vertex c splitting that complementary arc into lengths floor(ell/2) and
ceil(ell/2). Both lengths are at most floor(g/2), so each is a shortest
cycle distance between its endpoints. Therefore

\[d_C(a,b)+d_C(b,c)+d_C(c,a)=d+\ell=g.\]

The chosen c is distinct from a and b. This covers antipodal a,b on an
even cycle, as well as g=3: no endpoint or short-cycle exception is
needed.

Choose one sensor in every U\_w. This branch-resolving core has size L
and activates exactly the b branch-containing roots. Indeed every such
rooted tree contains a terminal fork, by taking a branching vertex
farthest from the root. If b\textgreater=2, complete two already active
roots to a geodesic triple using at most one additional cycle sensor. If
b=1, choose a second root and then a third completing the triple, adding
at most two sensors. If b=0, branch resolution is vacuous and any
geodesic triple of cycle vertices suffices. Apply the cited
geodesic-triple lemma in each case. QED.

For a metric basis, b\textgreater=2 now gives p\textless=1 by Lemma 6.
The same applies when b=1 and beta\textless=L+1. Lemma 7 leaves only
beta=L+2 in the one-branch-root case.

\subsection{7.4. The one-branch-active-root
case}\label{the-one-branch-active-root-case}

\textbf{Lemma 8 (one branch-active root).} If b=1, beta(G)=L+2, and B is
a metric basis, then p(G,B)\textless=1.

\emph{Proof.} In a finite rooted subcubic tree containing a branching
vertex, a deepest branching vertex has two terminal leafward threads;
therefore b=1 implies L\textgreater=1. All terminal forks lie in the
unique branch-containing rooted tree T\_r. Apply Lemma 6's exchange
construction to obtain a minimum fault-tolerant set F containing B\_0,
with one chosen B vertex in each U\_w. Let C=B without B\_0; then
\textbar C\textbar=2.

If F meets C, F union B has at most \textbar F\textbar+1 vertices and is
a sufficient superset. Otherwise F and C are disjoint. Since F resolves,
the imported characterization ensures at least two active cycle roots.
Choose z in F outside T\_r. Every B\_0 vertex lies inside T\_r, and no
vertex of C belongs to F, so z is not in B. Define S=B union (F without
z). Then \textbar S\textbar=\textbar F\textbar+1 and S contains B.

We check every possible sensor erasure from S. If a vertex of C is
erased, F without z remains, and resolves by fault tolerance of F. If a
vertex of F without B is erased, the resolving set B remains. It remains
to erase a vertex q of B\_0. The set F without \{z,q\} is
branch-resolving: each U\_w originally contained at least two vertices
of F, z lies outside T\_r, and deleting q removes at most one vertex
from one U\_w. Relative to the resolving set F without z, the root r
stays active because another F vertex remains in q's U\_w. Other active
roots are unchanged, as is occupancy of direct cycle threads. The
characterization implies that F without \{z,q\} resolves, so the larger
set S without q resolves. These cases exhaust S, proving fault
tolerance. Thus rho\textless=\textbar F\textbar+1 and p\textless=1. QED.

\subsection{7.5. No branch-active roots: odd
cycles}\label{no-branch-active-roots-odd-cycles}

Suppose b=0. Then each rooted attachment is a single path, possibly
consisting only of its cycle root, and L=0. Write a vertex on the path
attached at r as (r,t), where t\textgreater=0 is its distance from the
cycle. Its distance to a cycle vertex q is t+d\_C(r,q).

For an odd cycle of length 2h+1, h\textgreater=1, choose cycle vertices
a,b at distance h. Consider the root difference d\_C(r,a)-d\_C(r,b).
Along the h-edge arc from a to b its values are

\[-h,-h+2,\ldots,h.\]

At the h interior roots of the other arc, traversed from b to a, its
values are

\[h-1,h-3,\ldots,-h+1.\]

Each list has distinct entries, and their parities differ, so all 2h+1
roots have different differences. The two distances from (r,t) to a and
b identify r through their difference and then identify t from either
distance. Therefore \{a,b\} resolves every such graph. The case h=1
includes a triangle and gives differences -1,1,0, so no short-cycle
exception is suppressed.

No singleton can resolve a graph containing a cycle: if its sensor is in
an attached path at r (including r itself), the two cycle neighbors of r
have equal distance to that sensor. Hence beta=2. Lemma 5 gives
p\textless=1 for every metric basis in this odd-cycle case.

\subsection{7.6. No branch-active roots: even
cycles}\label{no-branch-active-roots-even-cycles}

Let the cycle have even length 2h with h\textgreater=2. Lemma 7 gives
beta\textless=3. A singleton cannot resolve, by the equal-distance
cycle-neighbor argument. Lemma 5 already handles beta=2, so assume
beta=3.

Every resolving set has at least two active cycle roots: if only r were
active, its two cycle neighbors would have equal distance to every
sensor. Furthermore a resolving set cannot have all its active roots in
a single antipodal pair \{a,a\emph{\}. The two cycle neighbors of a have
distance one to a and distance h-1 to a}, and thus equal distances to
every sensor in paths rooted at a or a*. They would remain unresolved.
Consequently a metric basis B has two non-antipodal active roots a,b.
Choose sensors x,y in B lying in their respective rooted paths.

Let a\emph{,b} be the cycle antipodes of a,b. The four roots
a,b,a\emph{,b} are distinct because a and b are distinct and
non-antipodal. The four sensors

\[F=\{x,y,a^*,b^*\}\]

also are distinct. After erasing any one of them, the three remaining
sensors occupy three distinct roots and contain an antipodal pair of
roots. Those three roots form a geodesic triple: for antipodes u,u* and
any other cycle vertex v, one of the two semicircles contains v, so
d\_C(u,v)+d\_C(v,u\emph{)=h; adding d\_C(u,u})=h gives the cycle length
2h. This reasoning includes a four-cycle, where any triple of distinct
roots contains an antipodal pair.

With b=0, the branch-resolving condition is vacuous. The cited
geodesic-triple lemma therefore makes every one-sensor deletion from F
resolving, so F is fault-tolerant and beta2\textless=4. For any graph
here, a fault-tolerant set of size k leaves a resolving set of size k-1
after an erasure; hence beta2\textgreater=beta+1. Since beta=3, beta2=4.
Finally, F contains x and y from B. Adding the remaining third vertex of
B produces a fault-tolerant superset of size at most five. Thus
rho(G,B)\textless=5 and p(G,B)\textless=1.

\subsection{7.7. Completion and
sharpness}\label{completion-and-sharpness}

For b\textgreater=2 the exchange bound and Lemma 7 give p\textless=1.
For b=1 the same bound handles beta\textless=L+1 and Lemma 8 handles
beta=L+2, exhausting Lemma 7's possibilities. For b=0, the odd-cycle
argument gives beta=2; on an even cycle beta\textless=3, and the
preceding construction handles the only remaining value beta=3. These
cases exhaust connected subcubic unicyclic graphs. Nonnegativity was
established in Section 2, completing Theorem 4. Proposition 3 supplies
equality p=1, so the bound is sharp.

The completed proof has been separately logically reviewed. The bounded
novelty assessment and its limitations are stated in Section 1. Neither
the computation nor separate-agent review is described as human peer
review.

\section{8. Software, provenance and
reproduction}\label{software-provenance-and-reproduction}

The original computation used CPython 3.13.1, 64-bit, build
\texttt{tags/v3.13.1:0671451,\ Dec\ 3\ 2024,\ MSC\ v.1942}, on Windows
11 platform string \texttt{Windows-11-10.0.26100-SP0}. The host reported
ten logical CPUs; the primary runner restricted execution to one logical
CPU. The graph generator, exact solvers and verifier use only the Python
standard library, with no third-party mathematical package, GPU,
stochastic choice or numerical tolerance in the original pilot. Section
9 uses separately recorded seeded input sampling with exact arithmetic.

The primary worker ran on 18 October 2026 from 10:56:33 to 10:56:44 UTC.
It used 10.390625 CPU seconds and 10.904554 wall seconds, with peak
resident working set 41,967,616 bytes. The launch plus synthetic
resource tests used 12.859375 aggregate CPU seconds. Saved-output
revalidation used 8.25 CPU seconds. Including a conservative two-second
allowance for initial control setup and syntax checking, the original
scientific/preflight charge was 23.109375 CPU seconds.

The separate external audit used 5.125 CPU seconds and 5.372638 wall
seconds on its successful run, with peak resident working set 87,089,152
bytes. It conservatively charged under 30 CPU seconds including an
instrumentation rerun. Thus the recorded combined charge is under
53.109375 CPU seconds against the original 7,200-second pilot budget.
These timings exclude manuscript authoring and document rendering. The
initial primary control setup failed closed on a page-rounded memory-cap
equality check; the audit's initial run failed during memory reporting
due to a ctypes handle-width error. Both are documented instrumentation
events, not mathematical discrepancies.

The primary Windows runner uses nested Job Objects, single-CPU affinity,
committed-memory and active-process limits, process hard working-set
maxima, and wall watchdogs. The external audit instead measured memory
and used a cooperative CPU guard, with no hard memory quota; it must not
be described as inheriting the primary worker's enforcement. Precise
page-rounded primary cap readbacks were checked but not retained. All
observed memory use was well below the 4 GB envelope.

Original commands, run from the pilot directory, were:

\begin{Shaded}
\begin{Highlighting}[]
\NormalTok{python launch}\OperatorTok{.}\FunctionTok{py}
\NormalTok{python audit\_saved}\OperatorTok{.}\FunctionTok{py}
\end{Highlighting}
\end{Shaded}

The launcher refuses to overwrite its existing launch ledger or run
directory. Reproduction should use a fresh isolated copy of the recorded
sources and frozen preparation artifacts, preserving the delivered
evidence. The public packaging step did not execute those scientific
commands again. The public pilot audit runs with
\texttt{python\ pilot-audit/audit.py} from a fresh release copy. Public
audit path adaptations are documented in PROVENANCE.json; original
source hashes are retained. Reference outputs must be preserved before
any rerun.

The evidence consists of \texttt{graphs.jsonl}, \texttt{bases.jsonl},
\texttt{certificates.jsonl}, \texttt{coverage.json},
\texttt{manifest.json}, the saved-artifact audit, and the separate
external-audit files. The manifest records input hashes before and after
execution and output hashes. \texttt{PROVENANCE.json} alongside this
manuscript pins the public files used here. The delivered ZIP's package
manifest pins its contents, including the audit copies.

Key SHA-256 hashes:

\begin{Shaded}
\begin{Highlighting}[]
\NormalTok{generate.py  72b31a974050a9eaa8d275758932756706782c3d47d93b023d686412c20f49cb}
\NormalTok{solve.py     e06c139310458cb54a81d6648ac161068d499f697ae578811375191b93b3963f}
\NormalTok{verify.py    2927b07bd449dd21e28a220417c53d55e38515d868353ab86eaa8e2dd6c501cc}
\NormalTok{graphs.jsonl 62e60f9bb99846634ed134a1340e4d054fc33600676bf68f1271f1e82db94f17}
\NormalTok{bases.jsonl  a9fc7c35757ac5e680c06ef315ff1fa6e0660b1168924759bed8acc94703d18e}
\NormalTok{certificates.jsonl}
\NormalTok{11569ee3a65bd4b46d209278a7a67b3c991bf2c5ef45fe062e53bc96270cfff1}
\end{Highlighting}
\end{Shaded}

\section{9. Computational robustness
extension}\label{computational-robustness-extension}

The local protocol, case plan and source hashes were recorded before the
recorded extension run; the runner checked those hashes before solving.
These files provide a local freeze, not an independently timestamped
preregistration. Their mutable local timestamps cannot establish the
absence of earlier unrecorded analyses. The seed was 20261018; every
sampled input is retained.

The 71 original graph classes through eight vertices supplied the
baseline seeds. Three unequal positive-integer edge-length assignments
in \{1,2,3,4\} per seed yielded 213 weighted cases. Attaching a leaf to
the first eligible degree-three vertex yielded 28 unicyclic degree-four
cases. Adding the first eligible nonedge yielded 70 subcubic two-cycle
cases. These deterministic, label-dependent variants are finite stress
tests, not exhaustive coverage of the larger classes. All minimum
ordinary metric bases of each case were included.

{\def\LTcaptype{none} % do not increment counter
\begin{longtable}[]{@{}lrrrr@{}}
\toprule\noalign{}
Family & Cases & Minimum bases & p=0 & p=1 \\
\midrule\noalign{}
\endhead
\bottomrule\noalign{}
\endlastfoot
Unit-length baseline & 71 & 697 & 679 & 18 \\
Unequal integer lengths & 213 & 1,337 & 1,317 & 20 \\
Unicyclic, degree four & 28 & 278 & 198 & 80 \\
Subcubic, two cycles & 70 & 672 & 602 & 70 \\
Total & 382 & 2,984 & 2,796 & 188 \\
\end{longtable}
}

All 382 cases completed, with no stopping-rule skips or premium above
one, across 68,576 subset entries. Three two-cycle cases repeat labelled
edge lists; the 70 cases represent 67 distinct labelled graphs. No
isomorphism-distinct extension count is claimed. Passing these cases
does not extend Theorem 4 beyond its unweighted, subcubic, unicyclic
hypotheses.

\subsection{9.1. Bounded-error
separation}\label{bounded-error-separation}

Each surviving coordinate has reading d(v,s)+e\_s, where
\textbar e\_s\textbar{} is at most epsilon; errors vary independently
and adversarially within closed intervals. One identified sensor erasure
deletes a coordinate, not a target vertex, and leaves graph distances
unchanged. For each target pair, sort the absolute coordinate gaps
\textbar d(u,s)-d(v,s)\textbar. Let m1(S) be the minimum over pairs of
the largest gap, and m2(S) the minimum over pairs of the second-largest
gap; missing gaps are zero. Uniform distinguishability without loss
requires m1(S)\textgreater2 epsilon; after any one identified erasure it
requires m2(S)\textgreater2 epsilon. This is an exact error-box
criterion, not a numerical tolerance rule. At equality the boxes
intersect, so an ambiguous permitted observation exists; this does not
mean every observation is ambiguous.

For any connected unit-edge graph with at least two vertices and any
exact fault-tolerant layout S, m2(S)=1. Integer distances and two
distinguishers imply m2(S)\textgreater=1. For adjacent targets u,v, the
triangle inequality gives \textbar d(u,s)-d(v,s)\textbar\textless=1 for
every sensor, implying m2(S)\textless=1. Thus its worst-case one-erasure
guarantee holds exactly for epsilon\textless1/2. The unit-edge
measurements are consistency checks of this generic fact, not a new
experimental or graph-class-specific discovery. Integrality likewise
makes every integer-distance FT layout pass at epsilon=1/4. An ordinary
minimum basis has m2=0 by minimality, since deleting any sensor cannot
leave a resolving set. No half-unit assertion is made for arbitrary
integer-weighted graphs.

The experiment evaluated epsilon=0,1/4,1/2,1 for each minimum basis, its
smallest-mask optimum FT augmentation, and one smallest-mask
unconstrained FT optimum per case. Among 1,337 weighted augmentation
rows, m2 was 1 for 1,300 rows, 2 for 24, and 3 for 13. Accordingly, 37
passed at epsilon=1/2 and 13 at epsilon=1. These are deterministic
counts conditioned on bases, possibly sharing layouts, not independent
trials or probabilities. Margin-optimal choices among tied layouts and
minimum robust sensor sets were not searched. Weighted margins depend on
length units and were not normalised across graph scales; larger
absolute margins do not establish improved physical sensing. No
stochastic error distribution or physical calibration is asserted.

\subsection{9.2. Independent verification and
evidence}\label{independent-verification-and-evidence}

A separately authored executable audit reproduced the complete finite
case plan, all optima and noise certificates; this is AI computational
review, not human peer review or physical validation. It independently
rebuilt the cases, used Bellman-Ford distances, and checked exhaustive
signature/deletion tables. Every one of the 4,675 per-case unique-layout
boundary certificates was verified by enumerating target pairs and
erased coordinates. Each contains an ambiguous pair and midpoint
readings at epsilon=m2/2. The reviewed files and exact hashes are
recorded in the bundled \texttt{extension-audit/AUDIT.md} and
\texttt{extension-audit/audit-result.json}; \texttt{extension\_audit.py}
has documented portable path adaptations; the original reviewed source
hash is retained. The reviewed result and noise-certificate hashes are:

\begin{Shaded}
\begin{Highlighting}[]
\NormalTok{results.jsonl}
\NormalTok{c202b7e2d20554255bf54f4ffdff6b069593321c3c8dff1ee22f25f6016d4c32}
\NormalTok{noise{-}boundary{-}certificates.jsonl}
\NormalTok{f3fe7a989b6e89d989005383934ee789c94009d2766ca29324576bef24fd4423}
\end{Highlighting}
\end{Shaded}

The producer's core and postcheck recorded 2.625 CPU seconds, with peak
working set below 52 MB. The separate audit recorded 0.90625 CPU seconds
and 54,587,392 bytes peak working set; minor inspection, setup,
summarisation and packaging overhead is not aggregated in these
subtotals. The core enforced a one-process, one-CPU Job Object with a
1,100-second wall watchdog and 3.3 GB committed-memory cap, plus a 3.0
GB working-set cap. The separate audit used a cooperative CPU guard and
measured memory, without a hard memory quota. Resource enforcement was
not stress-tested again by that audit. No new scientific cases were run
for manuscript integration.

\section{10. Interpretation and
limitations}\label{interpretation-and-limitations}

The quantity p isolates the cost of preserving a given metric basis in a
cardinality objective. For the example, retaining two original sensors
requires two additions and four total sensors; a freely redesigned
fault-tolerant solution uses three total sensors. This does not assign
monetary costs to adding, moving or removing a sensor, and it does not
establish an engineering advantage for either policy.

The main theorem assumes exact graph distances, equal edge lengths and
identification of the erased coordinate. Section 9 separately studies
finite weighted/topology variants and an uncalibrated bounded-error
model. It is not a model of water pressure, hydraulic flow, noisy
observations, leaks, unknown sensor errors or physical network
performance. No physical or network utility follows from the finite
result.

The computation excludes graphs of order above ten. Theorem 4 covers
arbitrary finite orders through a separate argument, not extrapolation.
Both are restricted to the stated graph class and fault model. The
finite stress tests in Section 9 do not establish a universal extension
to larger maximum degree, weighted or multicyclic graphs; directed
graphs, multigraphs and graph-vertex deletion faults remain outside the
analysis. Independent executable agreement still trusts Python, the
operating environment and the completeness/canonicalization arguments.
No proof assistant was used. The proof begins with an optimal
fault-tolerant set and is not by itself a polynomial-time optimization
algorithm. Its sensor-count interpretation does not demonstrate
noisy-distance robustness, hydraulic validity or field performance.

\section{11. AI assistance and draft
status}\label{ai-assistance-and-draft-status}

This draft was prepared under timofy's direction with AI assistance in
developing and checking mathematical arguments, implementing and
auditing computations, literature searching and comparison, and
manuscript drafting and editing. Separate AI sessions performed logical
and computational checks. Those checks are not independent human peer
review, journal review or a formal proof-assistant verification. No
affiliation or ORCID is asserted. This is a publicly shared AI-assisted
research draft. No human peer review has been undertaken and no journal
acceptance is claimed. The bounded search does not guarantee
originality; timofy coordinates the project; the recorded reviews are AI
logical and computational checks.

\section{References}\label{references}

{[}1{]} C. Hernando, M. Mora, P. J. Slater and D. R. Wood.
Fault-Tolerant Metric Dimension of Graphs. Convexity in Discrete
Structures, RMS Lecture Notes Series 5 (2008), 81-85.
HMSW-FaultTolerant.pdf

{[}2{]} J. Sedlar and R. Skrekovski. Vertex and edge metric dimensions
of unicyclic graphs. Discrete Applied Mathematics 314 (2022), 81-92.
DOI: 10.1016/j.dam.2022.02.022. Proof numbering: arXiv:2104.00577v1.
2104.00577

{[}3{]} M. Knor, J. Sedlar and R. Skrekovski. Fault Tolerance of Metric
Basis Can Be Expensive. Mediterranean Journal of Mathematics (2025).
DOI: 10.1007/s00009-025-02894-3. Checked preprint: arXiv:2410.03656v1.
2410.03656v1

{[}4{]} J. Geneson and S.-F. Tsai. Fault tolerance for metric dimension
and its variants. arXiv:2502.02731v1 (2025). Checked preprint version;
no numerical constant from Theorem 2.7 is used. 2502.02731v1

{[}5{]} A. Hakanen, V. Junnila, T. Laihonen and I. G. Yero. On Vertices
Contained in All or in No Metric Basis. Discrete Applied Mathematics 319
(2022), 407-423. 2103.08911

{[}6{]} B. Bergougnoux, O. Defrain and F. Mc Inerney. Enumerating
Minimal Solution Sets for Metric Graph Problems. Algorithmica 87 (2025),
712-735. Checked preprint: arXiv:2309.17419v2. 2309.17419v2

{[}7{]} A. Estrada-Moreno. The k-Metric Dimension of a Unicyclic Graph.
Mathematics 9 (2021), 2789. DOI: 10.3390/math9212789. 2789

{[}8{]} S. Prabhu, S. Klavzar, K. Bharani Dharan and S. Radha. Computing
fault-tolerant metric dimension of graphs using their primary subgraphs.
arXiv:2509.14300v2 (2026). 2509.14300v2

\end{document}
